\section{Behavioral diagnostics and archived research trajectories}
\label{sec:si-behavioral}

\subsection{Study materials and local inference settings}

We examined Cloudra's recorded research trajectories and constructed local questions about the molecules, geometries, and calculations in them. The trajectories retain the original model statements, tool inputs and outputs, and subsequent actions. The local tests vary which information the model receives and the order of its answer options. The supplementary data contain every question, expected answer, actual answer, request configuration, and source identifier.

The local deployment was archived as \texttt{qwen3.8:27b}, using a GGUF Q4\_K\_M quantization with a reported parameter size of 27.3B. The deployment digest is
\begin{quote}\small\ttfamily
22130167c4c20e20c7b71454612966ca8e8171e9b3cc8ab6ce8aa6cbfec79643.
\end{quote}
The local API reports a maximum context of 262,144 tokens. The short diagnostic requests use an 8,192-token context. Temperature is 0, seed 42, top-$k$ 0, top-$p$ 1, minimum-$p$ 0, repeat penalty 1, presence penalty 0, and draft prediction count 0. The initial identity and evidence-binding tests use the saved explicit ChatML rendering, an empty thinking block, thinking disabled, and a maximum of eight generated tokens; top-20 token log probabilities are retained. Native-chat controls use separate system and user messages, the model's native rendering, and a maximum of 1,024 generated tokens. The thinking setting is recorded for each request. A final choice is valid only if it exactly matches a permitted uppercase letter. The letter and its meaning are recorded separately.

\begin{table}[!htbp]
\centering\small
\caption{Experimental units and crossed conditions of the local diagnostic studies.}
\label{tab:si-diagnostic-design}
\begin{tabular}{p{0.24\linewidth}p{0.57\linewidth}r}
\toprule
Study & Frozen inputs and experimental factors & Calls\\
\midrule
Chemical identity & 16 pairs $\times$ 2 representations $\times$ 2 answer orders & 64\\
Distance judgment & 4 geometries $\times$ 2 information conditions $\times$ 2 answer orders $\times$ 2 thinking settings & 32\\
Target-angle feedback & 2 information conditions; 3 cyclic next-action permutations and 2 sufficiency permutations per condition & 10\\
Evidence associations & 4 chains $\times$ 2 association conditions $\times$ 2 presentations $\times$ 2 answer orders & 32\\
Native diagnostic follow-up & 20 identity inputs and 16 association inputs, each under 2 thinking settings & 72\\
Outstanding requirements & 3 state questions $\times$ 2 context conditions $\times$ 2 answer orders & 12\\
\bottomrule
\end{tabular}
\par\smallskip\begin{minipage}{\linewidth}\footnotesize
Each call is counted once. Representation, answer-order, and information variants use the same source molecule, geometry, or calculation chain. The single expanded-budget distance follow-up is additional to the 32-call study. The activation replacement experiment is described separately below.
\end{minipage}
\end{table}

\subsection{Chemical identity under equivalent representations and answer reversal}

Sixteen molecular pairs were derived from four task systems. They include pairs representing the same molecule, different connectivities with the same formula, changes in connectivity and formal-charge location, different double-bond stereochemistry, and protonation changes. Each pair was presented in its original notation and an equivalent rewriting, with both answer orders. In the forward order, A means the same chemical identity and B means different identities; in the swapped order the meanings exchange. The questions instruct the model to compare connectivity, bond orders, formal charges, and specified stereochemistry, allowing differences in traversal order, ring digits, and disconnected-component order.

Across the resulting 64 judgments, 58 answers were correct. For the same pair and answer order, equivalent rewriting changed the chemical judgment in four of 32 comparisons. For the same pair and representation, reversing the options changed the judgment in six of 32 comparisons. All six selected A in both option orders. In the same-formula pair \texttt{CC1CC=C1} and \texttt{C=C1CCC1}, the double bond lies within and outside the ring, respectively. The model answered A under all four representation/order conditions, even though A changed from ``same'' to ``different'' when the options were reversed.

Identical strings test whether the model recognizes a representation as matching itself. Equivalent rewritings test recognition of the same molecule in another notation. Different-identity controls test whether it detects changed connectivity or stereochemistry despite a shared molecular formula. All strings and expected judgments are saved with the diagnostic data.

\subsection{XYZ coordinates and explicitly measured distances}

Four NB1 geometries were selected from the archived research trajectories. The local question asks which of two specified atom pairs has the shorter Euclidean distance. The original-XYZ condition supplies the complete geometry in its saved atom order. The measured-distance condition supplies program-calculated values for the same atom pairs. Each condition appears in both answer orders and in direct-answer and thinking modes, yielding 32 native chat calls. Table~\ref{tab:si-distance-stimuli} gives the measured values used to construct the paired conditions; all atom indices are one-based.

\begin{table}[!htbp]
\centering\small
\caption{Distance comparisons derived from the four archived NB1 geometries.}
\label{tab:si-distance-stimuli}
\begin{tabular}{llrlr}
\toprule
Source workflow & Pair 1 & $d_1$ (\AA) & Pair 2 & $d_2$ (\AA)\\
\midrule
Guided & C2--C4 & 2.029 & O3--H8 & 1.248\\
Brief & C4--H11 & 2.197 & O7--H11 & 0.981\\
Brief & C4--H11 & 1.388 & O7--H11 & 1.362\\
Brief & C4--H11 & 2.386 & O7--H11 & 0.968\\
\bottomrule
\end{tabular}
\end{table}

\begin{table}[!htbp]
\centering\small
\caption{Complete outcomes of the native distance-judgment controls.}
\label{tab:si-distance-outcomes}
\begin{tabular}{llrrrr}
\toprule
Thinking & Information & Calls & Correct & Valid wrong & Budget-truncated\\
\midrule
Off & XYZ coordinates & 8 & 4 & 4 & 0\\
Off & Measured distances & 8 & 8 & 0 & 0\\
On & XYZ coordinates & 8 & 7 & 0 & 1\\
On & Measured distances & 8 & 8 & 0 & 0\\
\bottomrule
\end{tabular}
\par\smallskip\begin{minipage}{\linewidth}\footnotesize
A budget-truncated generation reached the 1,024-token ceiling while still awaiting a final choice. It remains in the eight-input denominator. A separate 2,048-token follow-up of that one input produced the correct choice after 1,353 generated tokens.
\end{minipage}
\end{table}

All eight direct-answer XYZ responses selected A, giving four correct answers; reversing the options changed the judgment for every source geometry. Measured distances gave eight correct answers in each thinking setting, with the same judgment under reversal. With thinking enabled, median generated-token counts over the eight paired inputs were 807.5 for XYZ and 86 for measured distances, including thinking tokens. Program-calculated distances let the model compare the requested quantities directly, as in program-aided language models.\citep{ref15}

\subsection{Target-specific geometric feedback in an E1 operation chain}

In the E1 trajectory, the agent seeks an anti-periplanar H7--C2--C3--Cl4 dihedral near $180^\circ$ for the transferring hydrogen H7. It rotates the C3-side fragment about C2--C3 by $+120^\circ$ and $-120^\circ$ to produce candidates A and B, and sets the C3--Cl4 distance to 2.35 \AA. The default report gives a dihedral involving H5. The agent requests the H7 dihedral needed to choose between the candidates.

\begin{table}[!htbp]
\centering\small
\caption{Dihedral measurements and circular residuals for the saved E1 candidates.}
\label{tab:si-target-dihedrals}
\begin{tabular}{llrr}
\toprule
Candidate & Atom quadruplet & Signed angle ($^\circ$) & Distance to $180^\circ$\\
\midrule
A & H5--C2--C3--Cl4 & $-55.7831$ & 124.2169\\
A & H7--C2--C3--Cl4 & 59.3752 & 120.6248\\
B & H5--C2--C3--Cl4 & 64.2169 & 115.7831\\
B & H7--C2--C3--Cl4 & 179.3752 & 0.6248\\
\bottomrule
\end{tabular}
\end{table}

The manuscript's Newman projections use the saved coordinates viewed along C2$\to$C3, with projected bond directions normalized for display. Table~\ref{tab:si-target-dihedrals} gives the signed angles and circular distances from the target. The agent subsequently submits candidate B as the transition-state guess, citing the H7 anti-periplanar arrangement and stretched C3--Cl4 bond.

Ten local tests ask what to do next and whether the available measurements suffice. The H5-only inputs withhold H7 values and XYZ coordinates; the H5-plus-H7 inputs add the measured H7 values. In all three cyclic option permutations, adding H7 changes the next-action answer from requesting that measurement to selecting candidate B. In two paired sufficiency tests, it changes the answer from requesting more measurements to judging the information sufficient. All values come from the saved geometries; only the measurements supplied to the model change.

\subsection{Single-token transfer of a reported numerical fact}

The activation experiment uses the archived local quantized model for a reported-distance comparison.\citep{ref16} We changed the first reported distance from 2.029 to 1.029~\AA. The second distance remains 1.248~\AA. This reverses the correct comparison. The clean and edited prompts each contain 124 tokens; their token IDs differ only at global position 53, corresponding to the initial digit of the changed measurement. The original-fact answer is B in the forward option order and A in the reversed order.

For each option order, the residual-stream state at position 53 is captured from the original prompt and transferred into the edited prompt at zero-based layers 7, 15, 31, and 47. The four layers are fixed in the experiment manifest. The readout is the full-vocabulary softmax probability mass assigned to the original-fact answer, summing its bare-letter and leading-space spellings. Table~\ref{tab:si-fact-patch} preserves the clean, edited, and transferred-state values.

\begin{table}[!htbp]
\centering\small
\caption{Probability of the original-fact answer under single-position activation replacement.}
\label{tab:si-fact-patch}
\begin{tabular}{lrr}
\toprule
Input / intervention & Forward: original answer B & Reversed: original answer A\\
\midrule
Original input & 0.954032 & 0.996932\\
Edited input & 0.008694 & 0.073626\\
Transfer at layer 7 & 0.029099 & 0.955908\\
Transfer at layer 15 & 0.015033 & 0.717660\\
Transfer at layer 31 & 0.007961 & 0.072117\\
Transfer at layer 47 & 0.008763 & 0.069810\\
\bottomrule
\end{tabular}
\par\smallskip\begin{minipage}{\linewidth}\footnotesize
The readout uses the full vocabulary, with bare and leading-space answer tokens combined. It is distinct from an A/B-renormalized probability or a proportion of correct task runs.
\end{minipage}
\end{table}

Two controls are executed in both option orders. Self-replacement transfers the edited prompt's own state at the changed token; the prefix control transfers at unchanged token position 0. Both preserve every vocabulary logit, with a maximum absolute difference of 0 from the edited baseline. With reversed options, transfers at layers 7 and 15 substantially increase the probability of the original-fact answer. With forward options, the changes are smaller and the edited answer remains preferred. Changing the state at the distance token therefore affects the answer in this test, with an effect that depends on option order.

\subsection{Evidence associations under fixed local calculation facts}

We used four calculation chains from three tasks to construct records with correct links and records with one changed link. The two NB1 chains come from the same run. Each record is presented as a flat list or with explicit links between calculations, in both answer orders, yielding $4\times2\times2\times2=32$ inputs. The changed link concerns the saddle-frequency geometry, IRC starting geometry, terminal-frame optimization input, or target species. Termination flags, convergence flags, energies, and imaginary-frequency values remain fixed.

The model is asked whether the calculations support the indicated reaction step. Linked inputs specify which geometry each frequency or IRC calculation uses, which terminal frame each endpoint optimization starts from, and which species the final minimum represents. The model answers all 16 linked inputs and 14 of 16 flat inputs correctly. Both flat-input errors occur with changed links and reversed options, in E1 and the second NB1 step. Explicit links state the calculation dependencies that the model must otherwise reconstruct from the list.\citep{ref17}

The native-chat follow-up uses the two chains on which these errors occurred, retaining every correct/changed-link, flat/linked, and option-order variant: 16 association questions per thinking mode. It also retains all four variants of five identity pairs with first-stage errors: 20 identity questions per mode. Table~\ref{tab:si-native-followup} reports all outcomes of these 72 calls. With thinking disabled, native chat gives 14/20 correct identity answers and 14/16 correct association answers, as in the initial tests. With thinking enabled, association answers reach 16/16; identity questions produce eight correct final choices and 12 generations truncated at the common 1,024-token limit.

\begin{table}[!htbp]
\centering\small
\caption{Complete outcomes for the selected native diagnostic follow-up.}
\label{tab:si-native-followup}
\begin{tabular}{llrrrr}
\toprule
Rendering / mode & Family & Calls & Correct & Valid wrong & Budget-truncated\\
\midrule
First-stage raw & Identity & 20 & 14 & 6 & 0\\
First-stage raw & Association & 16 & 14 & 2 & 0\\
Native, thinking off & Identity & 20 & 14 & 6 & 0\\
Native, thinking off & Association & 16 & 14 & 2 & 0\\
Native, thinking on & Identity & 20 & 8 & 0 & 12\\
Native, thinking on & Association & 16 & 16 & 0 & 0\\
\bottomrule
\end{tabular}
\par\smallskip\begin{minipage}{\linewidth}\footnotesize
The first-stage rows are the selected records from the existing 96-input identity/association study. The four native rows contain 72 additional calls. All saved input variants are included for each selected case; responses truncated by the output limit remain in the denominators.
\end{minipage}
\end{table}

\subsection{From a detected gap to the next research action}

The HA1 and NB1 trajectories show what the agents did after receiving diagnostics. In a Brief HA1 run, the agent repeatedly changed the source-minimum reference after compact feedback, while the IRC endpoint link remained unresolved. In a Guided NB1 run, the frequency calculation found zero imaginary modes and returned a normal-mode diagnostic. The Coordinator requested changes to the transition-state guess and optimization constraints; the next calculation used those changes. The saved sequences connect each diagnostic to the next instruction and calculation request.\citep{ref18}

In IA1, the delivery report has \texttt{complete=false}, a pending bidirectional short-IRC verification, and an empty formal-waiver list. After earlier positive comments about intermediate optimization, the Coordinator requests to stop with the study marked complete. The completion check returns the pending IRC requirement to the model. The next stop request marks the study partial.

The local IA1 test asks three questions: whether the latest report declares completion, whether a formal short-IRC waiver exists, and whether the missing calculation remains required. Each question is presented with the report alone or a quoted 12-message history, in both answer orders. All 12 answers match the recorded report fields and waiver list; all six option-reversed pairs preserve the same judgment.

\subsection{Discussion of the diagnostic scope}

The tests examine whether a model recognizes a molecule in different notations, compares distances, requests measurements for the specified atoms, uses a supplied numerical fact, detects changed calculation links, and identifies pending IRC verification. They describe the recorded local model and the tested inputs. The native follow-up uses cases selected after first-stage errors and retains all their variants; its counts apply to that selected set. The common output limit bounds response length; actual computation depends on prompt conditioning and the tokens generated in each thinking mode.

Activation replacement measures the effect of a changed state on the quantized Qwen model's distance-comparison answer. Establishing the original DeepSeek runs' internal mechanism or the chemical validity of their results requires separate evidence. In the E1 trajectory, Supervisor/workspace input intervenes between the measurements and the agent's use of candidate B. The separate local tests vary the supplied target measurements to examine their effect on the choice.

The completed IA1 tests use the report-only and 12-message inputs. A separate request with the full 306-message history contained 242,803 counted input tokens in a 262,144-token context. It reached the 3,600-s HTTP waiting limit while awaiting an answer, leaving the saved response empty. Of six planned full-history inputs, the first was attempted and the other five remain unattempted. This incomplete request is recorded separately from the 12 completed short-context judgments.
